\section{Zadanie 3}
\autor{Oskar Prusik}
Niech \(Z\) będzie zmienną losową o gęstości \(f(z) = \frac{1}{2} e^{-|z|}\) dla każdego \(z \in \mathbb{R}\).
Niech \(X\) i \(Y\) będą niezależnymi zmiennymi losowymi o rozkładzie wykładniczym z parametrem 1.
Policz funkcje tworzące momenty dla rozkładów \(Z\) i \(X - Y\).
Wywnioskuj, że te zmienne mają taki sam rozkład.

\subsection*{Rozwiązanie}

Zmienna \(Z\) ma gęstość \(f(z)=\frac12 e^{-|z|}\).
\[
	M_Z(t)=\ev{e^{tZ}}
	=\int_{-\infty}^{\infty} e^{tz}\cdot\frac12 e^{-|z|}\,dz
	=\frac12\int_{-\infty}^{0} e^{z(t+1)}\,dz
	+\frac12\int_{0}^{\infty} e^{z(t-1)}\,dz.
\]

Dla \(t\in(-1,1)\) mamy:
\[
	\frac12\int_{-\infty}^{0} e^{z(t+1)}\,dz = \frac12\cdot\frac{1}{t+1},\qquad
	\frac12\int_{0}^{\infty} e^{z(t-1)}\,dz = \frac12\cdot\frac{1}{1-t}.
\]

Zatem
\[
	M_Z(t)=\frac12\left(\frac{1}{1+t}+\frac{1}{1-t}\right)
	=\frac{1}{1-t^2},\qquad t\in(-1,1).
\]

\medskip

Niech \(X,Y\sim\mathrm{Exp}(1)\) będą niezależne. Wtedy ich MGF obliczamy:

\[
	M_X(t)=\ev{e^{tX}}
	=\int_{0}^{\infty} e^{tx} e^{-x}\,dx
	=\int_{0}^{\infty} e^{-(1-t)x}\,dx.
\]

Całka zbieżna dla \(1-t>0\), czyli \(t<1\):
\[
	M_X(t)=\left[ -\frac{e^{-(1-t)x}}{1-t} \right]_{0}^{\infty}
	=0-\left(-\frac{1}{1-t}\right)=\frac{1}{1-t},\qquad t<1.
\]

Analogicznie \(M_Y(t)=\frac{1}{1-t}\) dla \(t<1\).

Dla \(X-Y\):
\[
	M_{X-Y}(t)=\ev{e^{t(X-Y)}}
	=\ev{e^{tX}}\ev{e^{-tY}}
	=M_X(t)\,M_Y(-t)
	=\frac{1}{1-t}\cdot\frac{1}{1-(-t)}
	=\frac{1}{1-t^2},\qquad t\in(-1,1).
\]

\medskip
Jako, że funkcje tworzące są takie same i dobrze zdefiniowane w otoczeniu 0, to rozkład też jest taki sam.[Nie wiem jakiego formlizmu tu wymaga Gajdzica pisząc wywnioskuj  :) ]